\documentclass[border=10pt]{standalone}
\usepackage{tikz,ctex}
\usetikzlibrary{positioning}

\begin{document}

\begin{tikzpicture}[
    % 全局设置
    node distance = 1.6cm and 1.4cm,
    every node/.style = {
        draw,
        rounded corners,
        thick,
        align = center,
        font = \sffamily\small,
        minimum width = 2.0cm,
        minimum height = 0.9cm
    },
    % 样式定义
    centerStyle/.style = {
        fill = blue!45,
        text = white,
        font = \sffamily\bfseries\large,
        line width = 2pt,
        minimum size = 2.8cm,
        circle
    },
    branchTitle/.style = {
        font = \sffamily\bfseries,
        line width = 1.5pt
    },
    subNode/.style = {
        fill = white,
        font = \sffamily\small,
        minimum width = 2.2cm,
        minimum height = 0.8cm
    },
    leafNode/.style = {
        fill = white,
        font = \sffamily\footnotesize,
        minimum width = 1.9cm,
        minimum height = 0.7cm
    },
    arrowStyle/.style = {
        ->,
        thick,
        >=stealth,
        shorten >=2pt
    }
]

    % ================= 1. 中心节点 =================
    \node[centerStyle] (center) {6.4.误差函数 \\ Error Functions};

    % ================= 2. 上方分支：回归 (红) =================
    \node[branchTitle, fill=red!15, above=of center] (regression) {回归 \\ Regression};

    \node[subNode, above=of regression, xshift=-6cm] (gaussian) {高斯噪声假设 \\ Gaussian Noise (6.23)};
    \node[subNode, above=of regression, xshift=-2cm] (likelihood) {似然函数 \\ Likelihood (6.24)};
    \node[subNode, above=of regression, xshift=2cm] (sumsq) {平方和误差 \\ Sum-of-Squares (6.26)};
    \node[subNode, above=of regression, xshift=6cm] (noisevar) {噪声方差估计 \\ Noise Variance (6.27)};

    % ================= 3. 右侧分支：二分类 (蓝) =================
    \node[branchTitle, fill=orange!15, right=of center] (binary) {二分类 \\ Binary Classification};

    \node[subNode, right=of binary, yshift=4cm] (bernoulli) {伯努利分布 \\ Bernoulli (6.32)};
    \node[subNode, right=of binary, yshift=2cm] (sigmoidout) {Sigmoid 输出 \\ Logistic Sigmoid};
    \node[subNode, right=of binary, yshift=0cm] (crossentropy) {交叉熵误差 \\ Cross-Entropy (6.33)};
    \node[subNode, right=of binary, yshift=-2cm] (multibinary) {多独立二分类 \\ Multi-binary (6.35)};
    \node[subNode, right=of binary, yshift=-4cm] (labelnoise) {标签噪声 \\ Label Noise $\epsilon$};

    % ================= 4. 下方分支：多分类 (绿) =================
    \node[branchTitle, fill=green!15, below=of center] (multiclass) {多分类 \\ Multiclass Classification};

    \node[subNode, below=of multiclass, xshift=-8cm] (onehot) {1-of-K 编码 \\ 1-of-K Coding};
    \node[subNode, below=of multiclass, xshift=-4.5cm] (softmax) {Softmax 输出 \\ Softmax (6.37)};
    \node[subNode, below=of multiclass, xshift=0.5cm] (multicross) {多类交叉熵 \\ Multiclass Cross-Entropy (6.36)};
    \node[subNode, below=of multiclass, xshift=6cm] (degeneracy) {权重空间退化 \\ Weight-space Degeneracy};

    % ================= 5. 左侧分支：统一框架与导数 (紫) =================
    \node[branchTitle, fill=purple!15, left=of center] (unified) {统一框架与导数 \\ Unified Framework \& Derivatives};

    \node[subNode, left=of unified, yshift=4cm] (pairing) {误差函数与激活配对 \\ Error-Activation Pairing};
    \node[subNode, left=of unified, yshift=2cm] (identity) {回归：恒等 + 平方和 \\ Identity + Sum-of-Squares};
    \node[subNode, left=of unified, yshift=0cm] (logistic) {二分类：Sigmoid + 交叉熵 \\ Sigmoid + Cross-Entropy};
    \node[subNode, left=of unified, yshift=-2cm] (softmaxpair) {多分类：Softmax + 交叉熵 \\ Softmax + Cross-Entropy};
    \node[subNode, left=of unified, yshift=-4cm] (derivative) {统一导数形式 \\ $\partial E/\partial a_k = y_k - t_k$ (6.31)};

    % ================= 连线逻辑 =================
    % 中心 -> 四大分支标题
    \draw[arrowStyle, red] (center) -- (regression);
    \draw[arrowStyle, blue] (center) -- (binary);
    \draw[arrowStyle, green!60!black] (center) -- (multiclass);
    \draw[arrowStyle, purple] (center) -- (unified);

    % 分支标题 -> 子节点 (上方)
    \draw[arrowStyle, red!60] (regression) -- (gaussian.south);
    \draw[arrowStyle, red!60] (regression) -- (likelihood.south);
    \draw[arrowStyle, red!60] (regression) -- (sumsq.south);
    \draw[arrowStyle, red!60] (regression) -- (noisevar.south);

    % 分支标题 -> 子节点 (右侧)
    \draw[arrowStyle, blue!60] (binary) -- (bernoulli.west);
    \draw[arrowStyle, blue!60] (binary) -- (sigmoidout.west);
    \draw[arrowStyle, blue!60] (binary) -- (crossentropy.west);
    \draw[arrowStyle, blue!60] (binary) -- (multibinary.west);
    \draw[arrowStyle, blue!60] (binary) -- (labelnoise.west);

    % 分支标题 -> 子节点 (下方)
    \draw[arrowStyle, green!60!black] (multiclass) -- (onehot.north);
    \draw[arrowStyle, green!60!black] (multiclass) -- (softmax.north);
    \draw[arrowStyle, green!60!black] (multiclass) -- (multicross.north);
    \draw[arrowStyle, green!60!black] (multiclass) -- (degeneracy.north);

    % 分支标题 -> 子节点 (左侧)
    \draw[arrowStyle, purple!60] (unified) -- (pairing.east);
    \draw[arrowStyle, purple!60] (unified) -- (identity.east);
    \draw[arrowStyle, purple!60] (unified) -- (logistic.east);
    \draw[arrowStyle, purple!60] (unified) -- (softmaxpair.east);
    \draw[arrowStyle, purple!60] (unified) -- (derivative.east);

\end{tikzpicture}

\end{document}